\documentclass[a4paper]{article}
\usepackage{amsfonts}
\usepackage{amsmath}
\usepackage{amssymb}
\usepackage{graphicx}     
\usepackage{cancel}
\usepackage{color}  

\newcommand{\Q}{\mathbb{Q}}
\newcommand{\R}{\mathbb{R}}
\newcommand{\llim}{\lim\limits}
\newcommand{\dint}{\displaystyle\int}

\begin{document}
  
\noindent \large Auteur: Torben Petré \\
\noindent \large Studierichting: Informatica\\
\noindent \large Oefeningenreeks 6B nr. 21.4\\

\medskip

\normalsize

Bereken de $n$-de partieelsom van $\dfrac{2(n+2)}{n(n+1)(n+3)(n+4)}$ \\

$= \dfrac{A}{n}+\dfrac{B}{n+1}+\dfrac{C}{n+3}+\dfrac{D}{n+4}$ \\

$A = \dfrac{2(0+2)}{(0+1)(0+3)(0+4)} = \dfrac{1}{3}$ \\

$B = \dfrac{2(-1+2)}{-1(-1+3)(-1+4)} = -\dfrac{1}{3}$ \\

$C = \dfrac{2(-3+2)}{-3(-3+1)(-3+4)} = -\dfrac{1}{3}$ \\

$D = \dfrac{2(-4+2)}{-4(-4+1)(-4+3)} = \dfrac{1}{3}$ \\

$= \dfrac{1}{3n} - \dfrac{1}{3(n+1)} - \dfrac{1}{3(n+3)} + \dfrac{1}{3(n+4)}$ \\

Als we dit dan invullen voor $n = 4$, dan zien we dat deze rij teleskoperend is. \\

$n = 1$\\

$\dfrac{1}{3} - \dfrac{1}{6} - \dfrac{1}{12} + \dfrac{1}{15}$\\

$n = 2$\\

$\dfrac{1}{6} - \dfrac{1}{9} - \dfrac{1}{15} + \dfrac{1}{18}$\\

$n = 3$\\

$\dfrac{1}{9} - \dfrac{1}{12} - \dfrac{1}{18} + \dfrac{1}{21}$\\

$n = 4$\\

$\dfrac{1}{12} - \dfrac{1}{15} - \dfrac{1}{21} + \dfrac{1}{24}$\\

Er blijven enkel 4 termen over: \\

$n = 1$\\

$\dfrac{1}{3} - \cancel{\dfrac{1}{6}} - \dfrac{1}{12} + \cancel{\dfrac{1}{15}}$\\

$n = 2$\\

$\cancel{\dfrac{1}{6}} - \cancel{\dfrac{1}{9}} - \cancel{\dfrac{1}{15}} + \cancel{\dfrac{1}{18}}$\\

$n = 3$\\

$\cancel{\dfrac{1}{9}} - \cancel{\dfrac{1}{12}} - \cancel{\dfrac{1}{18}} + \cancel{\dfrac{1}{21}}$\\

$n = 4$\\

$\cancel{\dfrac{1}{12}} - \dfrac{1}{15} - \cancel{\dfrac{1}{21}} + \dfrac{1}{24}$\\

$= \dfrac{1}{3} - \dfrac{1}{12} - \dfrac{1}{3(n + 1)} + \dfrac{1}{3(n + 4)}$ \\ 

$= \dfrac{1}{4} - \dfrac{1}{3(n + 1)} + \dfrac{1}{3(n + 4)}$ \\ 		

$= \dfrac{3(n+1)(n+4)-4(n+4)+4(n+1)}{12(n+1)(n+4)}$ \\

$= \dfrac{3(n^4+5n+4)-4n-16+4n+4}{12(n^2+5n+4)}$ \\

$= \dfrac{3n^4+15n\cancel{+12}\cancel{-4n}\cancel{-16}\cancel{+4n}\cancel{+4}}{12(n^2+5n+4)}$ \\

$= \dfrac{\cancel{3}(n^2+5n)}{\cancel{12}(n^2+5n+4)}$ \\

$= \dfrac{n^2+5n}{4n^2+20n+16}$ \\

\begin{thebibliography}{999}
%\bibitem{a} Referentie
\end{thebibliography}
\end{document} 